\documentclass[../../../include/open-logic-section]{subfiles}

\begin{document}
\olfileid{sth}{card-arithmetic}{simp}

\olsection{ساده‌سازی جمع و ضرب}

چنان‌که معلوم می‌شود، جمع و ضرب ترامتناهیِ کاردینالی
\emph{بسیار} ساده است. این امر از آن‌جا ناشی می‌شود که کاردینال‌ها
(برخی) اعداد ترتیبی‌اند، و ازاین‌رو خوش‌ترتیب‌اند و می‌توان آن‌ها را
به شیوه‌ای معین دست‌کاری کرد. بااین‌حال، نشان‌دادن این مطلب
\emph{چندان} آسان نیست. برای آغاز به تعریفی ظریف نیاز داریم:

\begin{defn}
یک \emph{ترتیب کانونی}، یعنی $\canonord$، را بر زوج‌های اعداد ترتیبی
چنین تعریف می‌کنیم: $\tuple{\alpha_1, \alpha_2} \canonord
\tuple{\beta_1, \beta_2}$ اگر و تنها اگر یکی از شرط‌های زیر برقرار باشد:
\begin{enumerate}
	\item $\max(\alpha_1, \alpha_2) < \max(\beta_1, \beta_2)$؛ یا
	\item $\max(\alpha_1, \alpha_2) = \max(\beta_1, \beta_2)$ و
	$\alpha_1 < \beta_1$؛ یا
	\item $\max(\alpha_1, \alpha_2) = \max(\beta_1, \beta_2)$ و
	$\alpha_1 = \beta_1$ و $\alpha_2 < \beta_2$.
\end{enumerate}
\end{defn}

\begin{lem}
برای هر عدد ترتیبی $\alpha$، $\tuple{\alpha \times \alpha, \canonord}$
یک خوش‌ترتیب است.
\end{lem}

\begin{proof}
آشکارا $\canonord$ بر $\alpha \times \alpha$ تام است. فرض کنید نه
$\tuple{\alpha_1, \alpha_2}$ در ترتیب $\canonord$ از
$\tuple{\beta_1, \beta_2}$ کوچک‌تر باشد و نه دومی از اولی. در این صورت
$\max(\alpha_1, \alpha_2) = \max(\beta_1, \beta_2)$ و
$\alpha_1 = \beta_1$ و $\alpha_2 = \beta_2$؛ پس
$\tuple{\alpha_1, \alpha_2} = \tuple{\beta_1, \beta_2}$.

برای نشان‌دادن خوش‌ترتیبی، بگذارید $X \subseteq \alpha\times\alpha$
ناتهی باشد. چون $\alpha$ عددی ترتیبی است، عددی چون $\delta$ کوچک‌ترین
عضوِ $\Setabs{\max(\gamma_1, \gamma_2)}{\tuple{\gamma_1,
\gamma_2} \in X}$ است. اکنون از
$\Setabs{\tuple{\gamma_1,\gamma_2} \in X}{\max(\gamma_1, \gamma_2) =
\delta}$ همهٔ زوج‌ها جز آن‌هایی را که مؤلفهٔ نخستِ کمینه دارند کنار
بگذارید؛ در میان زوج‌های باقی‌مانده، زوجی که مؤلفهٔ دومِ کمینه دارد،
عضوِ $\canonord$-کمینهٔ $X$ است.
\end{proof}
\noindent
اکنون ملاحظه‌ای بسیار کوچک و ساده:

\begin{prop}\ollabel{simplecardproduct}
اگر $\cardeq{\alpha}{\beta}$، آنگاه
$\cardeq{\alpha \times \alpha}{\beta \times \beta}$.
\end{prop}

\begin{proof}
کافی است $f \colon \alpha \to \beta$ نگاشت
$\tuple{\gamma_1, \gamma_2} \mapsto
\tuple{f(\gamma_1), f(\gamma_2)}$ را القا کند.
\end{proof}
\noindent
اکنون همهٔ این مطالب را در اثبات لمی اساسی به کار می‌گیریم:
\begin{lem}\ollabel{alphatimesalpha}
برای هر عدد ترتیبی نامتناهی $\alpha$،
$\cardeq{\alpha}{\alpha \times \alpha}$.
\end{lem}

\begin{proof}
برای برهان خلف، بگذارید $\alpha$ کوچک‌ترین عدد ترتیبی نامتناهی باشد که
این حکم برای آن نادرست است. از
\olref[sfr][siz][zigzag]{natsquaredenumerable} داریم
$\cardeq{\omega}{\omega\times\omega}$، پس $\omega \in \alpha$.
افزون‌براین، $\alpha$ یک کاردینال است. وگرنه، برای برهان خلف،
$\card{\alpha} \in \alpha$؛ در نتیجه، بنا بر فرض،
$\cardeq{\card{\alpha}}{\card{\alpha} \times \card{\alpha}}$، و بنا بر
تعریف $\cardeq{\card{\alpha}}{\alpha}$؛ پس با استفاده از
\olref{simplecardproduct} خواهیم داشت
$\cardeq{\alpha}{\alpha\times\alpha}$.

اکنون برای هر $\tuple{\gamma_1, \gamma_2} \in \alpha \times \alpha$،
قطعهٔ زیر را در نظر بگیرید:
\begin{align*}
	\text{Seg}(\gamma_1, \gamma_2) &= \Setabs{\tuple{\delta_1, \delta_2} \in \alpha \times \alpha}{\tuple{\delta_1, \delta_2} \canonord \tuple{\gamma_1, \gamma_2}}
\end{align*}
اگر $\gamma = \max(\gamma_1, \gamma_2)$، آنگاه
$\tuple{\gamma_1, \gamma_2} \canonord \tuple{\gamma+1, \gamma + 1}$.
اگر $\gamma$ متناهی باشد، این قطعه متناهی است و بنابراین کاردینالیتهٔ
آن از کاردینال نامتناهی $\alpha$ کوچک‌تر است. اگر $\gamma$ نامتناهی
باشد، داریم:
\begin{align*}
	\text{Seg}(\gamma_1, \gamma_2) &
	\precsim ((\gamma \ordplus 1)\ordtimes (\gamma \ordplus 1))\\
	&\approx (\gamma \ordtimes \gamma)
	\text{، بنا بر \olref[ord-arithmetic][using-addition]{ordinfinitycharacter} و
	\olref{simplecardproduct}}\\
	&\approx \gamma \text{، بنا بر فرض استقرا}\\
	& \prec \alpha\text{، زیرا $\alpha$ یک کاردینال است.}
\end{align*}
پس $\ordtype{\alpha\times \alpha, \canonord} \leq \alpha$ و در نتیجه
$\cardle{\alpha \times \alpha}{\alpha}$. از سوی دیگر، البته
$\cardle{\alpha}{\alpha \times \alpha}$؛ بنابراین نتیجه از قضیهٔ
شرودر--برنشتاین به دست می‌آید.
\end{proof}

سرانجام به نتیجهٔ ساده‌کنندهٔ خود می‌رسیم:

\begin{thm}\ollabel{cardplustimesmax}
اگر $\cardfont{a}, \cardfont{b}$ کاردینال‌هایی نامتناهی باشند، آنگاه:
\[
	\cardfont{a}
\cardtimes \cardfont{b} = \cardfont{a} \cardplus \cardfont{b} =
\text{max}(\cardfont{a}, \cardfont{b}).
\]
\end{thm}

\begin{proof}
بدون کاستن از کلیت، فرض کنید
$\cardfont{a} = \max(\cardfont{a}, \cardfont{b})$. آنگاه با به‌کارگیری
\olref{alphatimesalpha}،
$\cardfont{a}\cardtimes\cardfont{a} = \cardfont{a} \leq \cardfont{a}
\cardplus \cardfont{b} \leq \cardfont{a} \cardplus \cardfont{a} \leq
\cardfont{a} \cardtimes \cardfont{a}$.
\end{proof}\noindent
به همین ترتیب، اگر $\cardfont{a}$ نامتناهی باشد، اجتماعِ
$\cardfont{a}$ مجموعه که هر یک اندازه‌ای $\leq\cardfont{a}$ دارند،
اندازه‌ای $\leq\cardfont{a}$ دارد:

\begin{prop}\ollabel{kappaunionkappasize}
بگذارید $\cardfont{a}$ کاردینالی نامتناهی باشد. برای هر عدد ترتیبی
$\beta \in \cardfont{a}$، بگذارید $X_\beta$ مجموعه‌ای باشد که
$\card{X_\beta} \leq \cardfont{a}$. آنگاه
$\card{\bigcup_{\beta \in \cardfont{a}} X_\beta} \leq \cardfont{a}$.
\end{prop}

\begin{proof}
برای هر $\beta \in \cardfont{a}$، یک !!a{injection} $f_\beta \colon
X_\beta \to \cardfont{a}$ برگزینید.\footnote{این نگاشت‌ها چگونه
«برگزیده» می‌شوند؟ بنگرید به \olref[sth][choice][countablechoice]{sec}.}
یک !!a{injection} $g \colon
\bigcup_{\beta \in \cardfont{a}} X_\beta \to \cardfont{a} \times
\cardfont{a}$ را با $g(v) = \tuple{\beta, f_\beta(v)}$ تعریف کنید، که
در آن $v \in X_\beta$ و برای هیچ $\gamma \in \beta$ نداریم
$v \in X_\gamma$. اکنون بنا بر \olref{cardplustimesmax}،
$\bigcup_{\beta \in \cardfont{a}} X_\beta \preceq
\cardfont{a} \times \cardfont{a} \approx \cardfont{a}$.
\end{proof}

\end{document}
